Notación: \(V\) = número de vértices, \(E\) = número de aristas, \(C\) = número de componentes conexas, \(\deg(v)\) = grado de \(v\). Se consideran grafos simples (sin aristas repetidas ni lazos).

\textbf{Grados y cantidad de aristas.}
\begin{props}
  \item \textbf{Lema del apretón de manos} (no dirigido): \(\sum_v \deg(v) = 2E\). En consecuencia, la cantidad de vértices de grado impar es siempre par.
  \item Máximo de aristas: \(\dfrac{V(V-1)}{2}\) en un grafo no dirigido y \(V(V-1)\) en uno dirigido.
  \item Un grafo conexo tiene \(E \ge V - 1\).
  \item \textbf{Número ciclomático:} \(E - V + C\) es la cantidad de ciclos independientes. Es 0 si y solo si el grafo no tiene ciclos.
\end{props}

\textbf{Árboles y bosques.}
\begin{props}
  \item Un grafo con \(V\) vértices es un árbol \(\iff\) es conexo y acíclico \(\iff\) es conexo con \(E = V-1\) \(\iff\) es acíclico con \(E = V-1\) \(\iff\) existe un único camino simple entre cada par de vértices.
  \item Un bosque con \(C\) componentes tiene \(E = V - C\).
  \item \textbf{Fórmula de Cayley:} existen \(V^{\,V-2}\) árboles etiquetados con \(V\) vértices (equivale a los árboles de expansión del grafo completo \(K_V\)).
\end{props}

\textbf{Grafos planares.}
\begin{props}
  \item \textbf{Teorema (fórmula) de Euler:} en un grafo planar conexo dibujado sin cruces, con \(F\) caras (contando la exterior), \(V - E + F = 2\). Con \(C\) componentes: \(V - E + F = 1 + C\).
  \item Si es planar simple con \(V \ge 3\): \(E \le 3V - 6\). Si además no tiene triángulos (por ejemplo, si es bipartito): \(E \le 2V - 4\).
\end{props}

\textbf{Otras propiedades.}
\begin{props}
  \item \textbf{Bipartito} \(\iff\) se puede 2-colorear \(\iff\) no tiene ciclos de longitud impar.
  \item \textbf{DAG:} un grafo dirigido admite orden topológico \(\iff\) no tiene ciclos.
  \item \textbf{Circuito euleriano} (recorre cada arista exactamente una vez y vuelve al inicio), en un grafo conexo no dirigido: existe \(\iff\) todos los vértices tienen grado par. Un \textbf{camino euleriano} (sin volver) existe \(\iff\) hay exactamente 0 o 2 vértices de grado impar; si son 2, el camino empieza en uno y termina en el otro.
  \item En un grafo dirigido, el circuito euleriano exige que el grado de entrada sea igual al grado de salida en todos los vértices (y que las aristas estén conectadas).
\end{props}
